\section*{Phụ lục A. Sai số theo từng horizon}

Bảng~\ref{tab:horizon} chi tiết hóa kết quả Test của Naive persistence và LightGBM
direct 7. Chạy lại pipeline với cấu hình trong mã nguồn và dữ liệu có mã băm bắt đầu
\texttt{\DatasetHashShort} sẽ tái tạo các giá trị này.

\begin{table}[H]
\centering
\caption{Sai số Test theo horizon (724 ngày gốc mỗi horizon)}
\label{tab:horizon}
\small
\begin{tabular}{crrrrrr}
\toprule
Horizon & \multicolumn{3}{c}{Naive persistence} & \multicolumn{3}{c}{LightGBM direct 7} \\
\cmidrule(lr){2-4}\cmidrule(lr){5-7}
 & RMSE & MAE & NSE & RMSE & MAE & NSE \\
\midrule
D+1 & 2{,}79 & 2{,}23 & 0{,}970 & 2{,}79 & 2{,}23 & 0{,}970 \\
D+2 & 3{,}93 & 3{,}12 & 0{,}940 & 3{,}93 & 3{,}12 & 0{,}940 \\
D+3 & 4{,}87 & 3{,}85 & 0{,}907 & 4{,}87 & 3{,}85 & 0{,}907 \\
D+4 & 5{,}60 & 4{,}42 & 0{,}877 & 5{,}60 & 4{,}43 & 0{,}877 \\
D+5 & 6{,}23 & 4{,}91 & 0{,}848 & 6{,}25 & 4{,}93 & 0{,}846 \\
D+6 & 6{,}68 & 5{,}16 & 0{,}824 & 6{,}72 & 5{,}19 & 0{,}822 \\
D+7 & 7{,}05 & 5{,}40 & 0{,}804 & 7{,}10 & 5{,}45 & 0{,}802 \\
\bottomrule
\end{tabular}

\end{table}

LightGBM chỉ nhỉnh hơn Naive ở D+1 (2,79 so với 2,79 m³/s sau làm tròn) và kém hơn
từ D+2 trở đi. Điều này cho thấy mô hình học máy chưa tạo được lợi thế ổn định so với
baseline persistence trên giai đoạn Test.

\section*{Phụ lục B. Hồ sơ số}

Mã nguồn, cấu hình, checkpoint, bảng kết quả và website tĩnh nằm trong
\texttt{submission.zip}; \texttt{source\_code.zip} chứa mã và test;
\texttt{latex\_overleaf.zip} chứa toàn bộ source LaTeX. Website chạy cục bộ theo
\texttt{README.md}; báo cáo không gắn URL công khai chưa được kiểm chứng. Run hiện tại:
\texttt{\RunId}; dữ liệu \texttt{\DatasetHashShort}; 14 test Python pass. PDF này được
tái tạo từ cùng nguồn và được kiểm tra còn 5 trang.
